{\small
\begin{threeparttable}
\begin{tabular}{llrrrr}
\toprule
Cohort & Period & Workers in the year & One dismissal & Final cohort & Months \\
\midrule
2014 & Pre-recession & 51,665,225 & 10,986,505 & 10,219,768 & 108 \\
2015 & Recession & 50,051,402 & 10,822,667 & 10,063,904 & 96 \\
2016 & Recession & 47,037,159 & 9,728,957 & 9,066,692 & 84 \\
2019 & Pre-pandemic & 45,903,410 & 8,290,226 & 8,179,238 & 48 \\
2020 & Pandemic & 45,213,299 & 8,395,762 & 8,257,022 & 36 \\
2021 & Pandemic & 47,334,862 & 7,789,668 & 7,582,350 & 24 \\
\addlinespace[3pt]
Total &  &  & 56,013,785 & 53,368,974 &  \\
\bottomrule
\end{tabular}
\begin{tablenotes}[flushleft]\footnotesize
\item \textit{Note:} ``Workers in the year'' counts those with at least one private job in the year. ``One dismissal'' counts those with a single dismissal without cause. The final cohort excludes workers with another private job active in the month of separation and those without a valid occupation code. ``Months'' is the follow-up through December 2023 for workers dismissed in December of the cohort year, the shortest in each cohort. Source: RAIS, authors' calculations.
\end{tablenotes}
\end{threeparttable}
}
